\section{Control Integration}

\subsection{Target Role of the Learned Model}

The long-term target is to use the learned predictor as a short-horizon transition model inside a controller, a lightweight simulator, or a model-based reinforcement-learning loop.
The intended interface is therefore not merely ``predict the next logged trajectory,'' but
\begin{equation}
\hat{\mathbf{x}}_{k+1:k+H}
=
f_{\theta}\!\left(\hat{\mathbf{x}}_{k}, \mathbf{u}_{k:k+H-1}, \mathbf{c}_{k}\right),
\end{equation}
possibly together with uncertainty and reliability information.

\subsection{Why Offline Screening Is Still Necessary}

At the current project stage, the repository should not claim verified closed-loop readiness.
Before the predictor is inserted into a control loop, it must first satisfy offline criteria related to numerical stability and transition fidelity.

The current offline metrics therefore act as a screening layer:
\begin{itemize}
  \item global error terms such as RMSE and MAE,
  \item final-step error,
  \item rollout-growth measures,
  \item tail-horizon error statistics,
  \item worst-component bias and non-finite-event counts.
\end{itemize}

These metrics do not prove closed-loop stability, but they do provide evidence about whether the learned model can serve as a plausible short-horizon solver candidate.

\subsection{Controller-Oriented Objective}

If the predictor is later used inside a finite-horizon controller, a generic objective can be written as
\begin{equation}
\min_{\mathbf{u}_{k:k+H-1}}
\sum_{j=0}^{H-1}
\left(
\left\|
\hat{\mathbf{x}}_{k+j} - \mathbf{x}^{\star}_{k+j}
\right\|_{\mathbf{Q}}^{2}
+
\left\|
\mathbf{u}_{k+j}
\right\|_{\mathbf{R}}^{2}
\right),
\end{equation}
subject to actuator and safety constraints.

In this context, the learned model is valuable only if its short-horizon rollout remains numerically stable and sensitive to the applied control sequence.

\subsection{Why the Single-Step Branch Matters}

The single-step transition branch is strategically important for control integration because it more directly matches the recursive computation required by control and simulation:
\begin{equation}
\hat{\mathbf{x}}_{k+1}
=
f_{\theta}^{\mathrm{step}}\!\left(\hat{\mathbf{x}}_{k}, \mathbf{u}_{k}, \mathbf{c}_{k}\right).
\end{equation}

Compared with a fixed future-block predictor, this form offers:
\begin{itemize}
  \item cleaner solver semantics,
  \item easier online recursion,
  \item clearer controller and RL interface design,
  \item more direct latency accounting in closed-loop evaluation.
\end{itemize}

\subsection{Current Boundary of Claims}

The present repository does \emph{not} yet establish:
\begin{itemize}
  \item closed-loop stability guarantees,
  \item high-fidelity physical-simulator equivalence,
  \item or deployment-grade safety certification.
\end{itemize}

What it does establish is a mathematically and engineeringly coherent bridge from offline sequence fitting to transition-solver evaluation:
\begin{itemize}
  \item causal proxy-state construction,
  \item transition-oriented training and selection,
  \item artifact-based offline screening,
  \item and a clear path toward controller-compatible interfaces.
\end{itemize}
